\section{Theory}

\subsection{Signal Model and Discretization}

In multi-b diffusion MRI the signal decay is a superposition of exponential decays from distinct diffusion compartments. One can either estimate the diffusivities $D_j$ and their fractions jointly---highly nonlinear and severely non-unique---or fix the diffusivities on a predefined grid and estimate only the fractions, reducing the problem to linear estimation at the cost of discretization. We adopt the latter: normalizing by the $b = 0$ signal, the discrete spectral model is
%
\begin{equation}
  \frac{S(b_i)}{S_0} = \sum_{j=1}^{K} R_j \exp(-b_i D_j) + \varepsilon_i,
  \quad i = 1, \ldots, M,
  \label{eq:signal_model}
\end{equation}
%
where $R_j \geq 0$ are the spectral fractions (volume fractions of each diffusion compartment), the $D_j$ are the $K$ diffusivity bin centers, and $\varepsilon_i \sim \mathcal{N}(0, \sigma^2)$ is additive Gaussian noise.
In matrix notation, $\svec / S_0 = \Umat \Rvec + \boldsymbol{\varepsilon}$, with design matrix $\Umat \in \mathbb{R}^{M \times K}$, $U_{ij} = \exp(-b_i D_j)$.
Here $M = 15$ b-values span $0$--$3500\;\text{s/mm}^2$ in steps of $250\;\text{s/mm}^2$, and $K = 8$ diffusivity bins lie at $D \in \{0.25, 0.5, 0.75, 1.0, 1.5, 2.0, 3.0, 20.0\}\;\umsmm$.
The design matrix is ill-conditioned ($\kappa(\Umat) \approx 2.8 \times 10^5$), reflecting the inherent difficulty of the inverse problem.

The choice of bins $D_j$ determines which spectral components can be resolved. Assuming known noise $\sigma$, the Fisher information matrix for the fractions is
%
\begin{equation}
  \mathbf{F} = \frac{1}{\sigma^2} \Umat^T \Umat,
  \quad F_{jk} = \frac{1}{\sigma^2} \sum_{i=1}^{M} \exp\!\big(-b_i(D_j + D_k)\big),
  \label{eq:fisher}
\end{equation}
%
\note{Sandy: please retrace Eq.~\ref{eq:fisher} and the $\Delta D \propto D^{3/2}$ derivation in the Appendix.}
and the Cram\'er--Rao lower bound (CRLB) gives the minimum variance of any unbiased estimator, $\text{Var}(\hat{R}_j) \geq [\mathbf{F}^{-1}]_{jj}$.
The off-diagonal structure of $\mathbf{F}$ shows why intermediate components are hard to separate: the Fisher correlation matrix $C_{jk} = F_{jk} / \sqrt{F_{jj} F_{kk}}$ exceeds $0.99$ for the components at $D = 0.50$, $0.75$, and $1.00\;\umsmm$ (Figure~\ref{fig:fisher}a), making them nearly indistinguishable from the available measurements. By contrast, $D = 0.25\;\umsmm$ contributes detectable signal across all 15 b-values, while fast components ($D \geq 3.0$) decay below the noise floor within the first few steps (Figure~\ref{fig:fisher}c). This correlation structure is independent of the fraction magnitudes, and is therefore the primary, mean-independent evidence for the identifiability limit reported below.

A two-compartment analysis makes the resolving limit explicit. For adjacent diffusivities $D$ and $D + \Delta D$ with fractions $p$ and $1-p$, in the regime $\Delta D \ll D$ and $b_{\max} D \gg 1$, the Fisher information for $p$ scales (Appendix) as $I \propto (\Delta D)^2 / D^3$, so equal precision across the spectrum would require a bin spacing growing as
%
\begin{equation}
  \Delta D \propto D^{3/2}.
  \label{eq:spacing_scaling}
\end{equation}
%
We use this only as a heuristic: our grid densifies toward slow diffusion ($\Delta D = 0.25\;\umsmm$ at the low end), but does not follow the rule strictly---we retain $D = 2.0\;\umsmm$, which strict spacing would omit, because it carries physiological (glandular) signal relevant to grading (Results). Numerical inversion of the full $8 \times 8$ Fisher matrix confirms this does not cost identifiability: intermediate components remain poorly conditioned regardless of bin placement, which only trades resolution between the two ends. The ill-posedness is intrinsic to the Laplace inversion, not an artifact of the grid.

The unconstrained CRLB greatly overstates the achievable error because it ignores the prior: for intermediate components it predicts standard deviations exceeding the unit range of a fraction (Figure~\ref{fig:fisher}b).
% TODO(Sandy): validate the van-Trees Bayesian-CRLB derivation plotted in Fig~\ref{fig:fisher}b.
The gap to the empirical posterior decomposes into two mechanisms rather than one ratio: incorporating the half-normal prior (the Bayesian, or van Trees, CRLB) tightens the unconstrained bound by up to two orders of magnitude, and the empirical NUTS posterior standard deviation is smaller still---by a further $2$--$74\times$ across components---because the non-negativity constraint and one-sided prior make the posterior non-Gaussian and truncated near the zero boundary, which the Gaussian-form bound does not capture (Figure~\ref{fig:fisher}b). \note{Sandy: confirm the van-Trees framing and the ``further $2$--$74\times$'' explanation; do we need a citation for the unconstrained and van-Trees CRLB?} The role of prior information is thus essential rather than cosmetic; the magnitude of this gain is revisited in the Discussion.

\subsection{MAP and Bayesian Estimation}

The maximum \emph{a posteriori} (MAP) estimate under a Gaussian (ridge) prior $R_j \sim \mathcal{N}(0, \lambda^{-1})$ combined with the physical non-negativity of compartment fractions solves the constrained quadratic program
%
\begin{equation}
  \hat{\Rvec}_{\text{MAP}} \;=\; \argmin_{\Rvec \,\in\, \mathbb{R}_{\geq 0}^{K}}
  \Big\{\; \big\lVert\, \svec/S_0 \;-\; \Umat\,\Rvec\, \big\rVert_2^{\,2}
  \;+\; \lambda\, \lVert \Rvec \rVert_2^{\,2}\, \Big\},
  \label{eq:map}
\end{equation}
%
where $\lambda > 0$ sets the regularization strength and $\mathbb{R}_{\geq 0}^{K}$ enforces non-negativity. For this ill-conditioned design the unconstrained Gaussian MAP $(\Umat^{T}\Umat + \lambda \mathbf{I})^{-1}\Umat^{T}(\svec/S_0)$ can lie outside the non-negative orthant, and projecting it onto $\mathbb{R}_{\geq 0}^{K}$ does not in general recover the constrained solution of (\ref{eq:map}); the difference is usually small but not negligible on a subset of ROIs. We therefore solve (\ref{eq:map}) directly as non-negative least squares on the augmented system $[\Umat;\, \sqrt{\lambda}\,\mathbf{I}]\,\Rvec = [\svec/S_0;\, \mathbf{0}]$ with an active-set solver. \note{Sandy: please confirm Eq.~\ref{eq:map} and the augmented-system notation.}

The MAP estimate is a point estimate: it neither quantifies how reliably each $R_j$ is determined nor accounts for the noise level $\sigma$, which silently rescales the penalty if mis-specified. A fully Bayesian treatment instead characterizes the joint posterior over $(\Rvec, \sigma)$, inferring spectrum and per-ROI noise together and making the ill-posedness explicit through the marginal posterior widths. We use the hierarchical model
%
\begin{align}
  R_j &\sim \text{HalfNormal}\!\left(\sigma_R = 1/\sqrt{\lambda}\right),
  \quad j = 1, \ldots, K \label{eq:prior_R} \\
  \sigma &\sim \text{HalfCauchy}(\beta = 0.01) \label{eq:prior_sigma} \\
  \svec / S_0 \mid \Rvec, \sigma &\sim \mathcal{N}(\Umat \Rvec,\; \sigma^2 \mathbf{I}).
  \label{eq:likelihood}
\end{align}
%
The half-normal prior on $\Rvec$ enforces non-negativity with smooth gradients for Hamiltonian Monte Carlo, and its variance $1/\lambda$ matches the ridge penalty, linking the MAP and Bayesian formulations; the half-Cauchy prior is a weakly informative choice \citep{gelman2006prior} that lets $\sigma$ be inferred jointly rather than fixed in advance \citep{sjolund2018bayesian}.
The posterior is sampled with the No-U-Turn Sampler (NUTS) \citep{hoffman2014nuts}, an adaptive variant of Hamiltonian Monte Carlo \citep{neal2011mcmc} that proposes joint moves across all components. A coordinate-wise Gibbs sampler has closed-form conditionals for this model (truncated-normal for $\Rvec$, conjugate inverse-gamma for $\sigma^2$); we use NUTS because the strong adjacent-component correlations (Figure~\ref{fig:fisher}a) make such coordinate updates mix slowly. \note{Sandy: I do not think I used the inverse-gamma $\sigma^2$ update in practice (I used the voxel-count SNR formula); confirm whether to include this or drop the Gibbs comparison.}
For each $R_j$, the posterior provides both a point estimate and a measure of identifiability through the coefficient of variation (CV $=$ posterior std $/$ posterior mean): components with CV $\ll 1$ are well constrained, those with CV near or above one effectively undetermined by the data.

\subsection{ADC as a Spectral Functional}

The apparent diffusion coefficient can be viewed as a functional of the spectrum. Given $\Rvec$, the synthesized signal is $\hat{S}(b) = \sum_j R_j \exp(-b D_j)$, and a monoexponential fit $\log \hat{S}(b) = \log \hat{S}_0 - b \cdot \text{ADC}$ over b-values below a cutoff yields $\text{ADC}(\Rvec)$. Log-domain least-squares fitting carries a known noise-distribution bias at low SNR \citep{kuczera2023truly} \stephan{this cites your earlier work---please confirm it is the right reference here} but is the clinical standard and is adopted for comparability. The mapping is nonlinear: a given fraction $R_j$ contributes to ADC differently depending on the overall spectral composition.

To identify which components ADC is most sensitive to, we define the sensitivity vector
%
\begin{equation}
  \frac{\partial \,\text{ADC}}{\partial R_j}\bigg|_{\Rvec = \bar{\Rvec}}
  \approx \frac{\text{ADC}(\bar{\Rvec} + \epsilon \, \mathbf{e}_j) - \text{ADC}(\bar{\Rvec})}{\epsilon},
  \label{eq:sensitivity}
\end{equation}
%
evaluated at a reference spectrum $\bar{\Rvec}$ (the average tumor or normal spectrum). We quantify how closely this sensitivity aligns with the optimal classification direction by the Pearson correlation between $\partial \text{ADC}/\partial \Rvec$ and the logistic-regression coefficient vector $\mathbf{w}$. The coefficients are standardized (each feature scaled to unit variance before fitting) so the comparison is on a common scale. Where this alignment holds, it reflects a concrete fact established in Results---both vectors concentrate on the two well-identified outer compartments---rather than any implicit optimality of the scalar itself.
